\subsection*{Context and prior work}
Arnold's fixed-point questions ask how much of the critical-point theory of a small autonomous Hamiltonian persists for a general Hamiltonian path \cite{Arnold1986}. Conley and Zehnder established the torus case by variational methods \cite{ConleyZehnder1983}. Floer's use of pseudoholomorphic curves turned the problem into a comparison between a Hamiltonian chain complex and the topology of the ambient manifold \cite{Floer1989}. Hofer--Salamon and Ono extended the theory in the weakly monotone setting \cite{HoferSalamon1995,Ono1995}; virtual-moduli methods subsequently yielded the general rational homological bound through work of Fukaya--Ono, Liu--Tian, and Ruan \cite{FukayaOno1999,LiuTian1998,Ruan1999}. The general rational bound counts Betti numbers. The Morse and stable Morse formulations ask for a geometric critical-point minimum, which can be larger.

The stable formulation is recorded explicitly in \cite[Conjecture~1.2]{Golovko2020}. Its distinction from the ordinary Morse formulation is essential: Damian exhibited closed manifolds with different ordinary and stable Morse numbers \cite{Damian2002}. In the positive direction, Dimitroglou Rizell and Golovko proved the stable Morse lower bound for generic Hamiltonians when both the symplectic area class and the first Chern class vanish on $\pi_2(M)$ \cite[Corollary~1.4]{DRG2017}. Their argument controls the simple homotopy type of the Floer complex. An independent proof was obtained by P\"oder Balkest\aa hl \cite[Corollary~1.2]{PoderBalkestahl2021}. The almost-quadratic convention for stable Morse functions used by Dimitroglou Rizell and Golovko agrees with ours by Lemma~\ref{lem:infinity-conventions}. Theorem~\ref{thm:main} disproves the unrestricted stable formulation, including its version that counts all fixed points.

Recent integral refinements explain why torsion alone does not force the bound in \eqref{eq:proposed-bound}. Abouzaid and Blumberg proved homological lower bounds over fields of arbitrary positive characteristic \cite{AbouzaidBlumberg2021}. Bai and Xu obtained an integral bound incorporating torsion at all primes \cite[Theorem~A]{BaiXu2022}, and their subsequent foundations give the same bound in \cite[Corollary~1.2]{BaiXu2026}. Its homology is graded modulo twice the minimal Chern number, so classes in different integer degrees may enter the same group. The geometric stable Morse number retains the integer degrees. For our examples the two bounds differ: Remark~\ref{rem:periodic-bound} computes the cyclically graded integral bound and shows that our fixed-point count exceeds it.

A separate construction gives nondegenerate Hamiltonian counterexamples to the ordinary Morse bound in real dimension twelve, with arbitrarily large deficit \cite[Theorem~1.1 and Corollary~1.2]{OpenAI2026Morse}. Those examples use a nontrivial fundamental group and the distinction between ordinary and stable Morse theory. In the construction below, the ambient manifold and the four fixed components of an auxiliary involution are simply connected, and Smale's theorem identifies both Morse numbers with the minimum allowed by integral homology \cite[Theorem~2.4]{Smale1963}. The separation instead comes from distributing torsion at two primes across different integer degrees and realizing the smaller Morse counts on the fixed components of an involution. The proof below is independent of that companion construction.
